% Options for packages loaded elsewhere
\PassOptionsToPackage{unicode}{hyperref}
\PassOptionsToPackage{hyphens}{url}
\PassOptionsToPackage{dvipsnames,svgnames,x11names}{xcolor}
%
\documentclass[
  11pt,
]{article}

\usepackage{amsmath,amssymb}
\usepackage{iftex}
\ifPDFTeX
  \usepackage[T1]{fontenc}
  \usepackage[utf8]{inputenc}
  \usepackage{textcomp} % provide euro and other symbols
\else % if luatex or xetex
  \usepackage{unicode-math}
  \defaultfontfeatures{Scale=MatchLowercase}
  \defaultfontfeatures[\rmfamily]{Ligatures=TeX,Scale=1}
\fi
\usepackage{lmodern}
\ifPDFTeX\else  
    % xetex/luatex font selection
    \setmainfont[]{Arial}
\fi
% Use upquote if available, for straight quotes in verbatim environments
\IfFileExists{upquote.sty}{\usepackage{upquote}}{}
\IfFileExists{microtype.sty}{% use microtype if available
  \usepackage[]{microtype}
  \UseMicrotypeSet[protrusion]{basicmath} % disable protrusion for tt fonts
}{}
\makeatletter
\@ifundefined{KOMAClassName}{% if non-KOMA class
  \IfFileExists{parskip.sty}{%
    \usepackage{parskip}
  }{% else
    \setlength{\parindent}{0pt}
    \setlength{\parskip}{6pt plus 2pt minus 1pt}}
}{% if KOMA class
  \KOMAoptions{parskip=half}}
\makeatother
\usepackage{xcolor}
\usepackage[a4paper, margin=0.8in]{geometry}
\setlength{\emergencystretch}{3em} % prevent overfull lines
\setcounter{secnumdepth}{5}
% Make \paragraph and \subparagraph free-standing
\makeatletter
\ifx\paragraph\undefined\else
  \let\oldparagraph\paragraph
  \renewcommand{\paragraph}{
    \@ifstar
      \xxxParagraphStar
      \xxxParagraphNoStar
  }
  \newcommand{\xxxParagraphStar}[1]{\oldparagraph*{#1}\mbox{}}
  \newcommand{\xxxParagraphNoStar}[1]{\oldparagraph{#1}\mbox{}}
\fi
\ifx\subparagraph\undefined\else
  \let\oldsubparagraph\subparagraph
  \renewcommand{\subparagraph}{
    \@ifstar
      \xxxSubParagraphStar
      \xxxSubParagraphNoStar
  }
  \newcommand{\xxxSubParagraphStar}[1]{\oldsubparagraph*{#1}\mbox{}}
  \newcommand{\xxxSubParagraphNoStar}[1]{\oldsubparagraph{#1}\mbox{}}
\fi
\makeatother


\providecommand{\tightlist}{%
  \setlength{\itemsep}{0pt}\setlength{\parskip}{0pt}}\usepackage{longtable,booktabs,array}
\usepackage{calc} % for calculating minipage widths
% Correct order of tables after \paragraph or \subparagraph
\usepackage{etoolbox}
\makeatletter
\patchcmd\longtable{\par}{\if@noskipsec\mbox{}\fi\par}{}{}
\makeatother
% Allow footnotes in longtable head/foot
\IfFileExists{footnotehyper.sty}{\usepackage{footnotehyper}}{\usepackage{footnote}}
\makesavenoteenv{longtable}
\usepackage{graphicx}
\makeatletter
\newsavebox\pandoc@box
\newcommand*\pandocbounded[1]{% scales image to fit in text height/width
  \sbox\pandoc@box{#1}%
  \Gscale@div\@tempa{\textheight}{\dimexpr\ht\pandoc@box+\dp\pandoc@box\relax}%
  \Gscale@div\@tempb{\linewidth}{\wd\pandoc@box}%
  \ifdim\@tempb\p@<\@tempa\p@\let\@tempa\@tempb\fi% select the smaller of both
  \ifdim\@tempa\p@<\p@\scalebox{\@tempa}{\usebox\pandoc@box}%
  \else\usebox{\pandoc@box}%
  \fi%
}
% Set default figure placement to htbp
\def\fps@figure{htbp}
\makeatother

\makeatletter
\@ifpackageloaded{caption}{}{\usepackage{caption}}
\AtBeginDocument{%
\ifdefined\contentsname
  \renewcommand*\contentsname{Table of contents}
\else
  \newcommand\contentsname{Table of contents}
\fi
\ifdefined\listfigurename
  \renewcommand*\listfigurename{List of Figures}
\else
  \newcommand\listfigurename{List of Figures}
\fi
\ifdefined\listtablename
  \renewcommand*\listtablename{List of Tables}
\else
  \newcommand\listtablename{List of Tables}
\fi
\ifdefined\figurename
  \renewcommand*\figurename{Figure}
\else
  \newcommand\figurename{Figure}
\fi
\ifdefined\tablename
  \renewcommand*\tablename{Table}
\else
  \newcommand\tablename{Table}
\fi
}
\@ifpackageloaded{float}{}{\usepackage{float}}
\floatstyle{ruled}
\@ifundefined{c@chapter}{\newfloat{codelisting}{h}{lop}}{\newfloat{codelisting}{h}{lop}[chapter]}
\floatname{codelisting}{Listing}
\newcommand*\listoflistings{\listof{codelisting}{List of Listings}}
\makeatother
\makeatletter
\makeatother
\makeatletter
\@ifpackageloaded{caption}{}{\usepackage{caption}}
\@ifpackageloaded{subcaption}{}{\usepackage{subcaption}}
\makeatother

\usepackage{bookmark}

\IfFileExists{xurl.sty}{\usepackage{xurl}}{} % add URL line breaks if available
\urlstyle{same} % disable monospaced font for URLs
\hypersetup{
  pdftitle={02\_GWAS\_Analysis},
  pdfauthor={Henry Janse van Rensburg},
  colorlinks=true,
  linkcolor={blue},
  filecolor={Maroon},
  citecolor={Blue},
  urlcolor={Blue},
  pdfcreator={LaTeX via pandoc}}


\title{02\_GWAS\_Analysis}
\author{Henry Janse van Rensburg}
\date{2026-11-16}

\begin{document}
\maketitle

\renewcommand*\contentsname{Table of contents}
{
\hypersetup{linkcolor=}
\setcounter{tocdepth}{2}
\tableofcontents
}

\section{Background}\label{background}

We conducted this experiment to investigate the genetic components of
BX-feedbacks. We have previously shown than BX-feedbacks are genetically
imprinted in Arabidopsis accessions. We used a GWAS approach to identify
candidate genes that are associated with BX-feedbacks. We grew 410
Arabidopsis accessions under controlled conditions on BXplus and BXminus
soils and measured the growth feedback. We used LIMIX and the
\texttt{gwaR} package to perform the GWAS analysis (@lippert2014) and
extract candidate genes (@hüther2020).

\section{Materials and methods}\label{materials-and-methods}

We grew 410 Arabidopsis accessions under controlled conditions. We used
a randomized block design with 3 replicates per accession. We grew the
plants in a growth chamber with automated phenotyping and extracted
phenotypes using ARADEEPOPSIS (@hüther2020). We measured the growth
feedback as the plant\_region\_area between plants grown on BXplus and
BXminus soils. We used the \texttt{gwaR} package to perform the GWAS
analysis.

\begin{itemize}
\tightlist
\item
  \textbf{Purpose}: To identify candidate genes that are associated with
  BX-feedbacks.
\end{itemize}

\section{Setup}\label{setup}

\subsubsection{Load Libraries and Install
packages}\label{load-libraries-and-install-packages}

\subsubsection{Import Phenotype table and SNP
Table}\label{import-phenotype-table-and-snp-table}

\section{Data Preparation}\label{data-preparation}

\subsection{Filter the SNP Table}\label{filter-the-snp-table}

We filter the SNP table and phenotype table to include only the data for
the first 20 days after germination (DAG 0 to DAG 20). A lot of
accessions started to flower at this point which affected the
phenotyping.

\begin{longtable}[]{@{}rrrrlrr@{}}
\caption{Table containing the SNP output data}\tabularnewline
\toprule\noalign{}
chrom & pos & log10\_p & pv & Significant & mac & dag \\
\midrule\noalign{}
\endfirsthead
\toprule\noalign{}
chrom & pos & log10\_p & pv & Significant & mac & dag \\
\midrule\noalign{}
\endhead
\bottomrule\noalign{}
\endlastfoot
3 & 21437549 & 8.990193 & 0 & Bonferroni & 5 & 12 \\
3 & 21437549 & 8.952353 & 0 & Bonferroni & 5 & 11 \\
3 & 21437549 & 8.870501 & 0 & Bonferroni & 5 & 13 \\
3 & 21437549 & 8.628093 & 0 & Bonferroni & 5 & 10 \\
3 & 21437549 & 8.473713 & 0 & Bonferroni & 5 & 14 \\
\end{longtable}

\subsection{Manhattan plot}\label{manhattan-plot}

\begin{itemize}
\tightlist
\item
  \textbf{Purpose}: The Manhattan plot visualizes the significance of
  SNPs across chromosomes, highlighting significant SNPs.
\end{itemize}

\subsection{Find Closest Genes to Top
SNPs}\label{find-closest-genes-to-top-snps}

We bioinformatically isolate the genes located closest to the top 10
SNPs identified in the GWAs analysis using the
\texttt{get\_nearest\_genes} function.

\begin{longtable}[]{@{}lrrlrrl@{}}
\caption{Table containing the SNPs and the genes located closest to each
of the top SNPs}\tabularnewline
\toprule\noalign{}
SNP\_rank & chrom & pos & Significant & log10\_p & mac & GeneId \\
\midrule\noalign{}
\endfirsthead
\toprule\noalign{}
SNP\_rank & chrom & pos & Significant & log10\_p & mac & GeneId \\
\midrule\noalign{}
\endhead
\bottomrule\noalign{}
\endlastfoot
1 & 3 & 21437549 & Bonferroni & 8.990193 & 5 & AT3G57890 \\
2 & 3 & 21437549 & Bonferroni & 8.952353 & 5 & AT3G57890 \\
3 & 3 & 21437549 & Bonferroni & 8.870501 & 5 & AT3G57890 \\
4 & 3 & 21437549 & Bonferroni & 8.628093 & 5 & AT3G57890 \\
5 & 3 & 21437549 & Bonferroni & 8.473713 & 5 & AT3G57890 \\
6 & 3 & 21437549 & Bonferroni & 8.199186 & 5 & AT3G57890 \\
7 & 3 & 16203882 & Bonferroni & 8.161670 & 12 & AT3G44630 \\
8 & 2 & 2259995 & Bonferroni & 8.125200 & 7 & AT2G05910 \\
9 & 3 & 16203882 & Bonferroni & 8.038630 & 11 & AT3G44630 \\
10 & 3 & 16203882 & Bonferroni & 7.859742 & 11 & AT3G44630 \\
\end{longtable}

\#Manhattan plot

\begin{figure}[H]

{\centering \pandocbounded{\includegraphics[keepaspectratio]{02_GWAS_Analysis_files/figure-pdf/Manhattat plot - manuscript-1.pdf}}

}

\caption{Figure 3a \textbar{} Genetic association mapping of growth
feedbacks in response to the BXplus microbiome. (a) Manhattan plot
showing the -log10-transformed P values for the top 500 single
nucleotide polymorphisms (SNPs) across the five chromosomes of
Arabidopsis, shown in alternating shades of grey. Associations were
tested for n = 410 accessions with a multi-trait linear mixed model
(limix; see Methods). SNPs that are significantly associated with growth
feedbacks after Bonferroni correction are shown in red. The biological
functions of genes adjacent to the three top-ranking SNPs are indicated.
(b) Average growth feedback among accessions carrying the alternative
allele of the SNP at position 16,203,882 (positive feedback) compared to
those with the reference allele (no or negative feedback). Each point is
one accession (mean growth feedback over three replicate pots per soil);
n = 11 accessions with the alternative allele and 399 with the reference
allele. Boxplots: centre line, median; box bounds, 25th and 75th
percentiles; whiskers, most extreme values within 1.5 × IQR
(interquartile range) of the box; all data points are shown. (c)
Feedback distribution curve with accessions carrying the alternative
allele of the SNP at position 16,203,882 indicated with black lines and
filled circles. (d) The genomic region (50 kb) flanking the SNP at
position 16,203,882 in the Col-0 reference genome (version TAIR10),
including the two TNL genes AT3G44630 (MMF1) and AT3G44670. SNP
significance is from the GWAS model in (a). (e) Venn diagram showing the
intersection between the genes located closest to the top-ranking 100
SNPs in (a), and genes differentially expressed (false discovery rate
\textless{} 0.05; empirical Bayes' shrinkage, DESeq2) in the root and
shoot transcriptome of Col-0 grown on the BXplus and BXminus microbiome
(data from ref (31)). Dataset 3.}

\end{figure}%

\section{Overlapping genes - RNAseq}\label{overlapping-genes---rnaseq}

\begin{figure}[H]

{\centering \pandocbounded{\includegraphics[keepaspectratio]{02_GWAS_Analysis_files/figure-pdf/VENN diagram-1.pdf}}

}

\caption{Figure 3e \textbar{} Genetic association mapping of growth
feedbacks in response to the BXplus microbiome. (a) Manhattan plot
showing the -log10-transformed P values for the top 500 single
nucleotide polymorphisms (SNPs) across the five chromosomes of
Arabidopsis, shown in alternating shades of grey. Associations were
tested for n = 410 accessions with a multi-trait linear mixed model
(limix; see Methods). SNPs that are significantly associated with growth
feedbacks after Bonferroni correction are shown in red. The biological
functions of genes adjacent to the three top-ranking SNPs are indicated.
(b) Average growth feedback among accessions carrying the alternative
allele of the SNP at position 16,203,882 (positive feedback) compared to
those with the reference allele (no or negative feedback). Each point is
one accession (mean growth feedback over three replicate pots per soil);
n = 11 accessions with the alternative allele and 399 with the reference
allele. Boxplots: centre line, median; box bounds, 25th and 75th
percentiles; whiskers, most extreme values within 1.5 × IQR
(interquartile range) of the box; all data points are shown. (c)
Feedback distribution curve with accessions carrying the alternative
allele of the SNP at position 16,203,882 indicated with black lines and
filled circles. (d) The genomic region (50 kb) flanking the SNP at
position 16,203,882 in the Col-0 reference genome (version TAIR10),
including the two TNL genes AT3G44630 (MMF1) and AT3G44670. SNP
significance is from the GWAS model in (a). (e) Venn diagram showing the
intersection between the genes located closest to the top-ranking 100
SNPs in (a), and genes differentially expressed (false discovery rate
\textless{} 0.05; empirical Bayes' shrinkage, DESeq2) in the root and
shoot transcriptome of Col-0 grown on the BXplus and BXminus microbiome
(data from ref (31)). Dataset 3.}

\end{figure}%

\begin{verbatim}
Overlap between Roots and GWAS:
\end{verbatim}

\begin{verbatim}
[1] "AT3G44630" "AT5G59220"
\end{verbatim}

\begin{verbatim}

Overlap between Shoots and GWAS:
\end{verbatim}

\begin{verbatim}
[1] "AT1G48100" "AT2G36080" "AT5G59220"
\end{verbatim}

\begin{itemize}
\tightlist
\item
  \textbf{Purpose}: Identifies the genes closest to the top 20 SNPs
  using the \texttt{get\_nearest\_genes} function.
\end{itemize}

\subsection{Plot the SNPs on Chromosome 3 Surrounding
SNPMMF1}\label{plot-the-snps-on-chromosome-3-surrounding-snpmmf1}

We zoom in on the top two genomic locations identified in the GWAS and
add the genes located closest to the SNPs.

\begin{itemize}
\tightlist
\item
  \textbf{Purpose}: The Manhattan plot for Chromosome 3 visualizes the
  significance of SNPs surrounding SNPMMF1, highlighting specific genes.
\end{itemize}

\section{Geographic distribution of
Accessions}\label{geographic-distribution-of-accessions}

We first create a world map that will be used to plot the geographic
distribution of Arabidopsis accessions based on their SNP response to
BX-feedbacks.

\subsection{Plot the Geographic Distribution of
Accessions}\label{plot-the-geographic-distribution-of-accessions}

We plot the distribution of accessions on the world map and color each
accession based on their BX-feedback response (either positive or
negative).

\begin{figure}

\centering{

\pandocbounded{\includegraphics[keepaspectratio]{02_GWAS_Analysis_files/figure-pdf/fig-SNPMMF1-Geogr-1.pdf}}

}

\caption{\label{fig-SNPMMF1-Geogr}Map showing the distribution of
accessions carrying SNPMMF1.}

\end{figure}%

\begin{itemize}
\tightlist
\item
  \textbf{Purpose}: The world map displays the geographic distribution
  of Arabidopsis accessions based on their SNP response to BX-feedbacks.
\end{itemize}

\section{Phenotypic effect of
SNPMMF1}\label{phenotypic-effect-of-snpmmf1}

We investigate the phenotypic effect of SNPMMF1 on the BX-feedback
response of Arabidopsis accessions. We use the
\texttt{phenotype\_by\_snp} function to extract the plant region area
and days after germination (DAG) data for SNPMMF1.

\begin{longtable}[]{@{}rlrrl@{}}
\caption{Table with plant region area by SNP.}\tabularnewline
\toprule\noalign{}
ACC\_ID & bx & plant\_region\_area & dag & SNP \\
\midrule\noalign{}
\endfirsthead
\toprule\noalign{}
ACC\_ID & bx & plant\_region\_area & dag & SNP \\
\midrule\noalign{}
\endhead
\bottomrule\noalign{}
\endlastfoot
9830 & + & 636 & 1 & FALSE \\
6092 & + & 413 & 1 & FALSE \\
6022 & + & 380 & 1 & FALSE \\
7424 & + & 568 & 1 & FALSE \\
9684 & - & 405 & 1 & FALSE \\
\end{longtable}

\subsection{Subset the Data to Phenotypes at Day
15}\label{subset-the-data-to-phenotypes-at-day-15}

We do this to focus on the phenotypic effects of SNPMM1 on growth. We
choose DAG15, since this is the latest time when the genetic association
of SNPMMF1 was still visible.

\begin{longtable}[]{@{}
  >{\raggedleft\arraybackslash}p{(\linewidth - 16\tabcolsep) * \real{0.0875}}
  >{\raggedright\arraybackslash}p{(\linewidth - 16\tabcolsep) * \real{0.0750}}
  >{\raggedleft\arraybackslash}p{(\linewidth - 16\tabcolsep) * \real{0.1250}}
  >{\raggedleft\arraybackslash}p{(\linewidth - 16\tabcolsep) * \real{0.1125}}
  >{\raggedleft\arraybackslash}p{(\linewidth - 16\tabcolsep) * \real{0.1250}}
  >{\raggedleft\arraybackslash}p{(\linewidth - 16\tabcolsep) * \real{0.1500}}
  >{\raggedright\arraybackslash}p{(\linewidth - 16\tabcolsep) * \real{0.1250}}
  >{\raggedright\arraybackslash}p{(\linewidth - 16\tabcolsep) * \real{0.1375}}
  >{\raggedleft\arraybackslash}p{(\linewidth - 16\tabcolsep) * \real{0.0625}}@{}}
\caption{Table with accession growth feedback and feedback
response}\tabularnewline
\toprule\noalign{}
\begin{minipage}[b]{\linewidth}\raggedleft
ACC\_ID
\end{minipage} & \begin{minipage}[b]{\linewidth}\raggedright
SNP
\end{minipage} & \begin{minipage}[b]{\linewidth}\raggedleft
-
\end{minipage} & \begin{minipage}[b]{\linewidth}\raggedleft
+
\end{minipage} & \begin{minipage}[b]{\linewidth}\raggedleft
area\_diff
\end{minipage} & \begin{minipage}[b]{\linewidth}\raggedleft
fold\_change
\end{minipage} & \begin{minipage}[b]{\linewidth}\raggedright
excessive
\end{minipage} & \begin{minipage}[b]{\linewidth}\raggedright
preference
\end{minipage} & \begin{minipage}[b]{\linewidth}\raggedleft
rank
\end{minipage} \\
\midrule\noalign{}
\endfirsthead
\toprule\noalign{}
\begin{minipage}[b]{\linewidth}\raggedleft
ACC\_ID
\end{minipage} & \begin{minipage}[b]{\linewidth}\raggedright
SNP
\end{minipage} & \begin{minipage}[b]{\linewidth}\raggedleft
-
\end{minipage} & \begin{minipage}[b]{\linewidth}\raggedleft
+
\end{minipage} & \begin{minipage}[b]{\linewidth}\raggedleft
area\_diff
\end{minipage} & \begin{minipage}[b]{\linewidth}\raggedleft
fold\_change
\end{minipage} & \begin{minipage}[b]{\linewidth}\raggedright
excessive
\end{minipage} & \begin{minipage}[b]{\linewidth}\raggedright
preference
\end{minipage} & \begin{minipage}[b]{\linewidth}\raggedleft
rank
\end{minipage} \\
\midrule\noalign{}
\endhead
\bottomrule\noalign{}
\endlastfoot
9851 & TRUE & 13541.167 & 59993.50 & 46452.33 & 4.430453 & plus & plus &
1 \\
7350 & TRUE & 12014.167 & 43003.83 & 30989.67 & 3.579427 & plus & plus &
2 \\
5921 & FALSE & 7656.833 & 36693.33 & 29036.50 & 4.792234 & plus & plus &
3 \\
5867 & FALSE & 11557.667 & 38961.17 & 27403.50 & 3.371024 & plus & plus
& 4 \\
9972 & TRUE & 27742.667 & 53904.83 & 26162.17 & 1.943030 & plus & plus &
5 \\
\end{longtable}

\begin{itemize}
\tightlist
\item
  \textbf{Purpose}: The data is filtered to include only plants that are
  exactly 15 days old. The mean plant region area for each ACC\_ID
  within each bx group is calculated. The data is reshaped so that bx
  treatments are in separate columns. The difference in plant region
  area between bx+ and bx- treatments is calculated, and the results are
  classified. The accessions are ranked based on the area difference,
  and the final data frame is saved to a CSV file.
\end{itemize}

\subsection{Plot the Relative Growth Feedback distribution of Accessions
with
SNPMMF1}\label{plot-the-relative-growth-feedback-distribution-of-accessions-with-snpmmf1}

\begin{figure}

\centering{

\pandocbounded{\includegraphics[keepaspectratio]{02_GWAS_Analysis_files/figure-pdf/fig-Growth-Feedback-SNPMMF1-1.pdf}}

}

\caption{\label{fig-Growth-Feedback-SNPMMF1}Figure 3c \textbar{} Genetic
association mapping of growth feedbacks in response to the BXplus
microbiome. (a) Manhattan plot showing the -log10-transformed P values
for the top 500 single nucleotide polymorphisms (SNPs) across the five
chromosomes of Arabidopsis, shown in alternating shades of grey.
Associations were tested for n = 410 accessions with a multi-trait
linear mixed model (limix; see Methods). SNPs that are significantly
associated with growth feedbacks after Bonferroni correction are shown
in red. The biological functions of genes adjacent to the three
top-ranking SNPs are indicated. (b) Average growth feedback among
accessions carrying the alternative allele of the SNP at position
16,203,882 (positive feedback) compared to those with the reference
allele (no or negative feedback). Each point is one accession (mean
growth feedback over three replicate pots per soil); n = 11 accessions
with the alternative allele and 399 with the reference allele. Boxplots:
centre line, median; box bounds, 25th and 75th percentiles; whiskers,
most extreme values within 1.5 × IQR (interquartile range) of the box;
all data points are shown. (c) Feedback distribution curve with
accessions carrying the alternative allele of the SNP at position
16,203,882 indicated with black lines and filled circles. (d) The
genomic region (50 kb) flanking the SNP at position 16,203,882 in the
Col-0 reference genome (version TAIR10), including the two TNL genes
AT3G44630 (MMF1) and AT3G44670. SNP significance is from the GWAS model
in (a). (e) Venn diagram showing the intersection between the genes
located closest to the top-ranking 100 SNPs in (a), and genes
differentially expressed (false discovery rate \textless{} 0.05;
empirical Bayes' shrinkage, DESeq2) in the root and shoot transcriptome
of Col-0 grown on the BXplus and BXminus microbiome (data from ref
(31)). Dataset 3.}

\end{figure}%

\begin{itemize}
\tightlist
\item
  \textbf{Purpose}: The bar plot visualizes the relative growth feedback
  for all accessions, highlighting specific accessions with SNPMMF1.
\end{itemize}

\subsection{Boxplot illustrating SNP Effect on growth
feedback}\label{boxplot-illustrating-snp-effect-on-growth-feedback}

\begin{figure}

\centering{

\pandocbounded{\includegraphics[keepaspectratio]{02_GWAS_Analysis_files/figure-pdf/fig-SNPMMF-growth-feedback-1.pdf}}

}

\caption{\label{fig-SNPMMF-growth-feedback}Figure 3b \textbar{} Genetic
association mapping of growth feedbacks in response to the BXplus
microbiome. (a) Manhattan plot showing the -log10-transformed P values
for the top 500 single nucleotide polymorphisms (SNPs) across the five
chromosomes of Arabidopsis, shown in alternating shades of grey.
Associations were tested for n = 410 accessions with a multi-trait
linear mixed model (limix; see Methods). SNPs that are significantly
associated with growth feedbacks after Bonferroni correction are shown
in red. The biological functions of genes adjacent to the three
top-ranking SNPs are indicated. (b) Average growth feedback among
accessions carrying the alternative allele of the SNP at position
16,203,882 (positive feedback) compared to those with the reference
allele (no or negative feedback). Each point is one accession (mean
growth feedback over three replicate pots per soil); n = 11 accessions
with the alternative allele and 399 with the reference allele. Boxplots:
centre line, median; box bounds, 25th and 75th percentiles; whiskers,
most extreme values within 1.5 × IQR (interquartile range) of the box;
all data points are shown. (c) Feedback distribution curve with
accessions carrying the alternative allele of the SNP at position
16,203,882 indicated with black lines and filled circles. (d) The
genomic region (50 kb) flanking the SNP at position 16,203,882 in the
Col-0 reference genome (version TAIR10), including the two TNL genes
AT3G44630 (MMF1) and AT3G44670. SNP significance is from the GWAS model
in (a). (e) Venn diagram showing the intersection between the genes
located closest to the top-ranking 100 SNPs in (a), and genes
differentially expressed (false discovery rate \textless{} 0.05;
empirical Bayes' shrinkage, DESeq2) in the root and shoot transcriptome
of Col-0 grown on the BXplus and BXminus microbiome (data from ref
(31)). Dataset 3.}

\end{figure}%

\begin{itemize}
\tightlist
\item
  \textbf{Purpose}: The boxplot with jitter visualizes the effect of
  SNPMMF1 on the growth feedback response.
\end{itemize}

\subsection{Combined SNPMMF1 effect
plots}\label{combined-snpmmf1-effect-plots}

\begin{itemize}
\tightlist
\item
  \textbf{Purpose}: The relative growth feedback plot and the SNP effect
  plot are combined into a single plot grid.
\end{itemize}

\section{\texorpdfstring{\textbf{5. Effect of SNPs linked to
SNPMMF1}}{5. Effect of SNPs linked to SNPMMF1}}\label{effect-of-snps-linked-to-snpmmf1}

Here we identify SNPs that are linked to SNPMMF1 using the
\texttt{linkage\_phenotypes} function and calculate the mean plant
region area at each SNP for acessions with and without the SNP.

\subsection{Identify linked SNPs}\label{identify-linked-snps}

\begin{itemize}
\tightlist
\item
  \textbf{Purpose}: The average plant region area at each SNP is
  calculated using the \texttt{linkage\_phenotypes} function.
\end{itemize}

\subsubsection{Calculate the Mean Plant Region Area at Each SNP for each
accession}\label{calculate-the-mean-plant-region-area-at-each-snp-for-each-accession}

\subsubsection{Plot the effect of this linked SNP inside MMF at position
16200156}\label{plot-the-effect-of-this-linked-snp-inside-mmf-at-position-16200156}

\subsubsection{Plot the effect of SNP at position
16218427}\label{plot-the-effect-of-snp-at-position-16218427}

\subsection{Calculate the Mean Plant Region Area at Each
SNP}\label{calculate-the-mean-plant-region-area-at-each-snp}

\begin{itemize}
\tightlist
\item
  \textbf{Purpose}: The mean plant region area at each SNP is
  calculated, and the data is reshaped to show bx treatments in separate
  columns. The difference in plant region area between bx+ and bx-
  treatments is calculated, and the data is sorted by position.
\end{itemize}

\section{Genomic exploration}\label{genomic-exploration}

\subsubsection{Load Arabidopsis Genome
Data}\label{load-arabidopsis-genome-data}

We load the Arabidopsis genome data from the Ensembl Plants database
using the \texttt{biomaRt} package.

\subsubsection{Format Arabidopsis Genome Data as
GRanges}\label{format-arabidopsis-genome-data-as-granges}

We format the Arabidopsis genome data as a GRanges object, which is
suitable for genomic range operations. We also prepare gene labels for
plotting by filtering and mutating the genomic data. Finally, we filter
the Arabidopsis genome data to include only the specific region
surrounding SNPMMMF1 on chromosome 3.

\begin{longtable}[]{@{}
  >{\raggedright\arraybackslash}p{(\linewidth - 8\tabcolsep) * \real{0.2388}}
  >{\raggedright\arraybackslash}p{(\linewidth - 8\tabcolsep) * \real{0.2388}}
  >{\raggedleft\arraybackslash}p{(\linewidth - 8\tabcolsep) * \real{0.2239}}
  >{\raggedleft\arraybackslash}p{(\linewidth - 8\tabcolsep) * \real{0.1940}}
  >{\raggedleft\arraybackslash}p{(\linewidth - 8\tabcolsep) * \real{0.1045}}@{}}
\caption{Table with genomic region surrounding SNPMMF1.}\tabularnewline
\toprule\noalign{}
\begin{minipage}[b]{\linewidth}\raggedright
ensembl\_gene\_id
\end{minipage} & \begin{minipage}[b]{\linewidth}\raggedright
chromosome\_name
\end{minipage} & \begin{minipage}[b]{\linewidth}\raggedleft
start\_position
\end{minipage} & \begin{minipage}[b]{\linewidth}\raggedleft
end\_position
\end{minipage} & \begin{minipage}[b]{\linewidth}\raggedleft
strand
\end{minipage} \\
\midrule\noalign{}
\endfirsthead
\toprule\noalign{}
\begin{minipage}[b]{\linewidth}\raggedright
ensembl\_gene\_id
\end{minipage} & \begin{minipage}[b]{\linewidth}\raggedright
chromosome\_name
\end{minipage} & \begin{minipage}[b]{\linewidth}\raggedleft
start\_position
\end{minipage} & \begin{minipage}[b]{\linewidth}\raggedleft
end\_position
\end{minipage} & \begin{minipage}[b]{\linewidth}\raggedleft
strand
\end{minipage} \\
\midrule\noalign{}
\endhead
\bottomrule\noalign{}
\endlastfoot
AT3G44665 & 3 & 16216097 & 16216809 & 1 \\
AT3G44610 & 3 & 16188129 & 16192365 & -1 \\
AT3G44660 & 3 & 16214964 & 16215785 & -1 \\
AT3G44630 & 3 & 16195848 & 16201225 & 1 \\
AT3G44635 & 3 & 16199123 & 16200095 & -1 \\
\end{longtable}

The table of all the genes in the region of SNPMMF1.

\subsubsection{Plot the Genomic Region Surrounding SNPMMF1 and
phenotypes}\label{plot-the-genomic-region-surrounding-snpmmf1-and-phenotypes}

We plot the phenotypic differences across SNPs and overlay gene
annotations to visualize the relationship between SNPs, their location
relative to genes and their effect on phenotypes.

\begin{itemize}
\tightlist
\item
  \textbf{Purpose}: The plot visualizes the phenotype differences across
  SNPs and overlays gene annotations to show the relationship between
  SNPs and genes.
\end{itemize}

\begin{figure}[H]

{\centering \pandocbounded{\includegraphics[keepaspectratio]{02_GWAS_Analysis_files/figure-pdf/Genomic region urrounding top SNP-1.pdf}}

}

\caption{Figure 3d \textbar{} Genetic association mapping of growth
feedbacks in response to the BXplus microbiome. (a) Manhattan plot
showing the -log10-transformed P values for the top 500 single
nucleotide polymorphisms (SNPs) across the five chromosomes of
Arabidopsis, shown in alternating shades of grey. Associations were
tested for n = 410 accessions with a multi-trait linear mixed model
(limix; see Methods). SNPs that are significantly associated with growth
feedbacks after Bonferroni correction are shown in red. The biological
functions of genes adjacent to the three top-ranking SNPs are indicated.
(b) Average growth feedback among accessions carrying the alternative
allele of the SNP at position 16,203,882 (positive feedback) compared to
those with the reference allele (no or negative feedback). Each point is
one accession (mean growth feedback over three replicate pots per soil);
n = 11 accessions with the alternative allele and 399 with the reference
allele. Boxplots: centre line, median; box bounds, 25th and 75th
percentiles; whiskers, most extreme values within 1.5 × IQR
(interquartile range) of the box; all data points are shown. (c)
Feedback distribution curve with accessions carrying the alternative
allele of the SNP at position 16,203,882 indicated with black lines and
filled circles. (d) The genomic region (50 kb) flanking the SNP at
position 16,203,882 in the Col-0 reference genome (version TAIR10),
including the two TNL genes AT3G44630 (MMF1) and AT3G44670. SNP
significance is from the GWAS model in (a). (e) Venn diagram showing the
intersection between the genes located closest to the top-ranking 100
SNPs in (a), and genes differentially expressed (false discovery rate
\textless{} 0.05; empirical Bayes' shrinkage, DESeq2) in the root and
shoot transcriptome of Col-0 grown on the BXplus and BXminus microbiome
(data from ref (31)). Dataset 3.}

\end{figure}%

\section{Export data to Dataset3}\label{export-data-to-dataset3}

\section{Session Information}\label{session-information}

\begin{itemize}
\tightlist
\item
  \textbf{Purpose}: Print the session information, including details
  about the R version, loaded packages, and the system environment. This
  is useful for reproducibility, as it documents the computational
  environment in which the analysis was performed.
\end{itemize}

\begin{verbatim}
R version: R version 4.5.1 (2025-06-13) 
\end{verbatim}

\begin{verbatim}
Attached packages:
\end{verbatim}

\begin{verbatim}
openxlsx - 4.2.8.1 
ggnewscale - 0.5.2 
plyranges - 1.30.1 
GenomicRanges - 1.62.1 
Seqinfo - 1.0.0 
IRanges - 2.44.0 
S4Vectors - 0.48.1 
BiocGenerics - 0.56.0 
generics - 0.1.4 
biomaRt - 2.66.2 
fstcore - 0.10.0 
maps - 3.4.3 
ggvenn - 0.1.19 
readxl - 1.5.0 
knitr - 1.52 
cowplot - 1.2.0 
magrittr - 2.0.5 
gwaR - 0.6 
lubridate - 1.9.5 
forcats - 1.0.1 
stringr - 1.6.0 
dplyr - 1.2.1 
purrr - 1.2.2 
readr - 2.2.0 
tidyr - 1.3.2 
tibble - 3.3.1 
ggplot2 - 4.0.3 
tidyverse - 2.0.0 
\end{verbatim}




\end{document}
